\chapter{期望}

\section{关于测度的积分}

\subsection*{可测函数}

设 $(\Omega,\mathcal F,\mu)$ 为测度空间，即 $\Omega$ 为非空集合，$\mathcal F$ 为其上的 $\sigma$ 域，$\mu:\mathcal F\to[0,+\infty]$ 满足 $\mu(\varnothing)=0$，并对两两不交的 $A_1,A_2,\ldots\in\mathcal F$ 满足
\[
\mu\left(\bigcup_{n=1}^{\infty}A_n\right)=\sum_{n=1}^{\infty}\mu(A_n).
\]

\begin{mydefinition}{可测函数}{measurable-function}
若函数 $X:\Omega\to\mathbb R$ 对每个 Borel 集 $B\in\mathcal B(\mathbb R)$ 都满足
\[
X^{-1}(B)=\{\omega\in\Omega:X(\omega)\in B\}\in\mathcal F,
\]
则称 $X$ 为 $\mathcal F$-\textbf{可测函数}。当 $\mu$ 为概率测度时，这就是随机变量的定义。
\end{mydefinition}

\begin{example}[实直线上的可测函数]
若 $g:\mathbb R\to\mathbb R$ 是 Borel 可测函数，则对任意定义在 $\mathcal B(\mathbb R)$ 上的测度 $\mu$，$g$ 都是测度空间 $(\mathbb R,\mathcal B(\mathbb R),\mu)$ 上的可测函数。特别地，恒等函数及连续实函数都是可测函数。
\end{example}

\subsection*{可测函数的性质}

\begin{proposition}{可测性的半直线判别}{measurable-halfline}
对函数 $X:\Omega\to\mathbb R$，下列条件等价：
\begin{enumerate}
\item $X$ 可测；
\item 对每个 $t\in\mathbb R$，$\{X\leq t\}\in\mathcal F$；
\item 对每个 $t\in\mathbb R$，$\{X<t\}\in\mathcal F$。
\end{enumerate}
\end{proposition}

% 补足由半直线生成 Borel sigma 域的判别证明。
\begin{proof}
\begin{enumerate}
\item 由 (1) 推出 (2)。因为 $(-\infty,t]$ 为 Borel 集，其原像 $\{X\leq t\}$ 属于 $\mathcal F$。
\item 由 (2) 推出 (3)。对任意 $t\in\mathbb R$，
\[
\{X<t\}=\bigcup_{n=1}^{\infty}\{X\leq t-1/n\}\in\mathcal F.
\]
若 $X(\omega)<t$，可取充分大的 $n$ 使 $1/n\leq t-X(\omega)$；反向包含由 $t-1/n<t$ 得到。
\item 由 (3) 推出 (1)。首先，
\[
\{X\leq t\}=\bigcap_{n=1}^{\infty}\{X<t+1/n\}\in\mathcal F.
\]
若 $X(\omega)>t$，取 $1/n<X(\omega)-t$，便知该点不在右侧交集中，所以等式成立。

令
\[
\mathcal H=\{B\subseteq\mathbb R:X^{-1}(B)\in\mathcal F\}.
\]
原像运算保持全集、补集和可数并，故 $\mathcal H$ 为 $\sigma$ 域。上式说明它包含所有 $(-\infty,t]$。这些半直线生成 $\mathcal B(\mathbb R)$：它们的可数并给出开半直线，取差得到开区间，而每个开集都是有理端点开区间的可数并。因此 $\mathcal B(\mathbb R)\subseteq\mathcal H$，即 $X$ 可测。
\end{enumerate}
\end{proof}

\begin{proposition}{可测函数的运算}{measurable-operations}
设 $X,Y:\Omega\to\mathbb R$ 可测，则：
\begin{enumerate}
\item 对任意 $a\in\mathbb R$，$aX$ 可测；
\item $X+Y$ 可测；
\item $XY$ 可测；
\item 若处处有 $Y\ne0$，则 $X/Y$ 可测。
\end{enumerate}
\end{proposition}

\begin{proof}
\begin{enumerate}
\item 当 $a>0$ 时，
\[
\{aX\leq t\}=\{X\leq t/a\}\in\mathcal F.
\]
当 $a<0$ 时，不等号反向，
\[
\{aX\leq t\}=\{X\geq t/a\}=\Omega\setminus\{X<t/a\}\in\mathcal F.
\]
当 $a=0$ 时，该集合在 $t<0$ 时为空集，在 $t\geq0$ 时为 $\Omega$。故三种情形均可测。
\item 对任意 $t\in\mathbb R$，
\[
\{X+Y<t\}
=\bigcup_{q\in\mathbb Q}\bigl(\{X<q\}\cap\{Y<t-q\}\bigr).
\]
若 $X(\omega)+Y(\omega)<t$，由有理数的稠密性，可取 $X(\omega)<q<t-Y(\omega)$，故该点属于右侧。反之，右侧任一集合中的点均满足 $X+Y<q+(t-q)=t$。右侧为可数个可测集合的并，故 $X+Y$ 可测。
% 补证性质：源 PDF 省略乘积可测性的证明，此处补全。
\item 对任意实值可测函数 $Z$，当 $t<0$ 时 $\{Z^2\leq t\}=\varnothing$；当 $t\geq0$ 时，
\[
\{Z^2\leq t\}=\{Z\geq-\sqrt t\}\cap\{Z\leq\sqrt t\}\in\mathcal F.
\]
所以 $Z^2$ 可测。由前两项，$X-Y$ 可测，再利用
\[
XY=\frac{(X+Y)^2-(X-Y)^2}{4},
\]
得到 $XY$ 可测。
% 补证性质：源 PDF 省略商可测性的证明，此处补全。
\item 定义 $g(0)=0$，并对 $y\ne0$ 定义 $g(y)=1/y$。函数 $g$ 在 $(-\infty,0)$ 和 $(0,\infty)$ 上连续；任意开集的原像是这两个区间中的开集之并，再按是否包含 $g(0)$ 决定是否添上 $\{0\}$，故为 Borel 集。因此 $g$ 是 Borel 可测函数。

对任意 Borel 集 $B$，有
\[
\{g(Y)\in B\}=Y^{-1}(g^{-1}(B))\in\mathcal F,
\]
所以 $g(Y)$ 可测。由 $Y$ 处处非零，$g(Y)=1/Y$，再由上一项可知 $X/Y=Xg(Y)$ 可测。
\end{enumerate}
\end{proof}

\textbf{注：}\quad 实值函数列的上、下确界和上、下极限可能取 $+\infty$ 或 $-\infty$。记 $\overline{\mathbb R}=\mathbb R\cup\{-\infty,+\infty\}$；若 $Z:\Omega\to\overline{\mathbb R}$ 对每个实数 $t$ 都满足 $\{Z\leq t\}\in\mathcal F$，则称 $Z$ 为\textbf{扩展实值可测函数}。用 $\{Z<t\}$ 作判别与此等价，证明中的可数并、交公式仍成立。

\begin{proposition}{可数上下确界与极限的可测性}{measurable-limits}
设 $X_1,X_2,\ldots$ 为实值可测函数，则
\[
\inf_{n\geq1}X_n,\qquad \sup_{n\geq1}X_n,\qquad
\liminf_{n\to\infty}X_n,\qquad \limsup_{n\to\infty}X_n
\]
均为扩展实值可测函数。其中逐点定义
\[
\liminf_{n\to\infty}X_n=\sup_{n\geq1}\inf_{k\geq n}X_k,\qquad
\limsup_{n\to\infty}X_n=\inf_{n\geq1}\sup_{k\geq n}X_k.
\]
进一步，若 $X_n$ 逐点收敛于有限值函数 $X$，则 $X$ 为实值可测函数。
\end{proposition}

\begin{proof}
\begin{enumerate}
\item 对每个 $t\in\mathbb R$，
\[
\left\{\inf_{n\geq1}X_n<t\right\}
=\bigcup_{n=1}^{\infty}\{X_n<t\}\in\mathcal F.
\]
当下确界严格小于 $t$ 时，至少有一项严格小于 $t$；反向包含直接由下确界的定义得到。
\item 对每个 $t\in\mathbb R$，
\[
\left\{\sup_{n\geq1}X_n\leq t\right\}
=\bigcap_{n=1}^{\infty}\{X_n\leq t\}\in\mathcal F.
\]
因为上确界不超过 $t$ 当且仅当每一项均不超过 $t$。
\item 对每个 $a\in\mathbb R$，
\[
\left\{\liminf_{n\to\infty}X_n\leq a\right\}
=\bigcap_{m=1}^{\infty}\bigcap_{n=1}^{\infty}
 \bigcup_{k=n}^{\infty}\{X_k<a+1/m\}.
\]
若某点属于左侧，则每个尾部的下确界都不超过 $a$，严格小于 $a+1/m$，所以对任意 $m,n$，总有 $k\geq n$ 使 $X_k<a+1/m$，该点属于右侧。

反过来，若某点属于右侧，则对任意 $m,n$，有
\[
\inf_{k\geq n}X_k\leq a+1/m.
\]
先对 $n$ 取上确界，得 $\liminf X_n\leq a+1/m$，再令 $m\to\infty$，得 $\liminf X_n\leq a$。右侧只涉及可数并、交，因而属于 $\mathcal F$。
\item 由数列上、下极限的关系，
\[
\limsup_{n\to\infty}X_n=-\liminf_{n\to\infty}(-X_n).
\]
上一项适用于实值可测函数列 $-X_n$。若 $Z$ 为扩展实值可测函数，则
\[
\{-Z\leq t\}=\Omega\setminus\{Z<-t\}\in\mathcal F,
\]
故取负仍可测，上极限也可测。当逐点极限存在且有限时，它等于下极限，因而为实值可测函数。
\end{enumerate}
\end{proof}

\subsection*{简单函数及其逼近}

\begin{mydefinition}{简单函数}{simple-function}
若 $A_1,\ldots,A_N\in\mathcal F$ 两两不交，且 $\bigcup_{i=1}^N A_i=\Omega$，$a_1,\ldots,a_N$ 为实数，则函数
\[
X=\sum_{i=1}^N a_i I_{A_i}
\]
称为\textbf{简单函数}。其中 $I_A$ 为集合 $A$ 的示性函数，即在 $A$ 上取 $1$，在 $\Omega\setminus A$ 上取 $0$。
\end{mydefinition}

\textbf{注：}\quad 对 $A\in\mathcal F$，
\[
\{I_A\leq t\}=
\begin{cases}
\varnothing,&t<0,\\
\Omega\setminus A,&0\leq t<1,\\
\Omega,&t\geq1.
\end{cases}
\]
所以 $I_A$ 可测，由可测函数的有限次线性运算可知，每个简单函数都可测。

\begin{mytheorem}{非负可测函数的简单函数逼近}{simple-approximation}
若 $X:\Omega\to[0,+\infty]$ 可测，则存在非负简单函数列 $\{X_n\}$，使得
\[
0\leq X_n\leq X_{n+1}\leq X,\qquad
\lim_{n\to\infty}X_n(\omega)=X(\omega)\quad(\omega\in\Omega).
\]
\end{mytheorem}

% 沿用二进分割与高度截断构造，补足单调性和逐点收敛的验证。
\begin{proof}
对每个正整数 $n$，令
\[
X_n=\sum_{k=0}^{n2^n-1}\frac{k}{2^n}
 I_{\{k/2^n\leq X<(k+1)/2^n\}}
 +nI_{\{X\geq n\}}.
\]
这些集合均可测、两两不交，且覆盖 $\Omega$，所以 $X_n$ 是非负简单函数。该构造把 $[0,n)$ 分成长度为 $2^{-n}$ 的区间，在每个区间取左端点的值，而把不小于 $n$ 的部分截断为 $n$，故 $0\leq X_n\leq X$。

固定 $\omega$，若 $x=X(\omega)<+\infty$，则
\[
X_n(\omega)=\min\left\{n,\frac{\lfloor2^n x\rfloor}{2^n}\right\},
\]
其中 $\lfloor u\rfloor$ 表示不超过 $u$ 的最大整数。由
\[
\lfloor2^{n+1}x\rfloor\geq2\lfloor2^n x\rfloor
\]
可知二进向下取整的值随 $n$ 不减，截断高度 $n$ 也不减，所以 $X_{n+1}(\omega)\geq X_n(\omega)$。当 $n>x$ 时，截断不再发生，且
\[
0\leq x-X_n(\omega)<2^{-n},
\]
从而 $X_n(\omega)\to x$。

若 $X(\omega)=+\infty$，则每个 $n$ 都有 $X_n(\omega)=n$，同样单调趋于 $X(\omega)$。因此结论在每个 $\omega$ 上均成立。
\end{proof}

\subsection*{非负简单函数的积分}

\begin{mydefinition}{非负简单函数的积分}{integral-simple}
设 $X$ 为测度空间 $(\Omega,\mathcal F,\mu)$ 上的非负简单函数，写成
\[
X=\sum_{i=1}^N a_i I_{A_i},\qquad a_i\geq0,
\]
其中 $A_1,\ldots,A_N\in\mathcal F$ 两两不交且并为 $\Omega$，称为 $\Omega$ 的一个\textbf{有限可测划分}。定义
\[
\int_\Omega X(\omega)\,\mu(d\omega)
:=\sum_{i=1}^N a_i\mu(A_i),
\]
也简记为 $\int_\Omega X\,d\mu$。积分取值于 $[0,+\infty]$。
\end{mydefinition}

\textbf{注：}\quad 本节在非负扩展实数的乘法中约定 $0\cdot(+\infty)=0$。因此，即使 $\mu(A_i)=+\infty$，零系数项仍贡献 $0$；零函数的积分恒为 $0$。

\textbf{注：}\quad 上述积分不依赖于非负简单函数所采用的划分和表示，因而是良定义的。

% 积分良定义性的证明在材料中未展开；用两个有限划分的共同细化补全。
\begin{proof}
设同一非负简单函数具有两种表示
\[
X=\sum_{i=1}^N a_iI_{A_i}
 =\sum_{j=1}^M b_jI_{B_j},\qquad a_i,b_j\geq0,
\]
其中 $\{A_i\}$ 与 $\{B_j\}$ 均为 $\Omega$ 的有限可测划分。集合 $A_i\cap B_j$ 两两不交且并为 $\Omega$，构成这两个划分的共同细化。

在每个非空的 $A_i\cap B_j$ 上，$X$ 同时等于 $a_i$ 和 $b_j$，故 $a_i=b_j$；若交集为空，则其测度为 $0$，相应的两项均为 $0$。利用测度的有限可加性，
\begin{align*}
\sum_{i=1}^N a_i\mu(A_i)
&=\sum_{i=1}^N\sum_{j=1}^M a_i\mu(A_i\cap B_j)\\
&=\sum_{i=1}^N\sum_{j=1}^M b_j\mu(A_i\cap B_j)\\
&=\sum_{j=1}^M b_j\mu(B_j).
\end{align*}
所以两种表示给出的积分值相同。
\end{proof}

\begin{proposition}{非负简单函数积分的性质}{integral-simple-properties}
在测度空间 $(\Omega,\mathcal F,\mu)$ 上，有下列结论。
\begin{enumerate}
\item 对任意 $A\in\mathcal F$，
\[
\int_\Omega I_A\,d\mu=\mu(A).
\]
\item 若 $X$ 是非负简单函数，则 $\int_\Omega X\,d\mu\geq0$。
\item 若 $X$ 是非负简单函数，$a\geq0$ 为实常数，则
\[
\int_\Omega aX\,d\mu=a\int_\Omega X\,d\mu.
\]
\item 若 $X,Y$ 都是非负简单函数，则
\[
\int_\Omega(X+Y)\,d\mu
=\int_\Omega X\,d\mu+\int_\Omega Y\,d\mu.
\]
\item 若非负简单函数 $X,Y$ 处处满足 $X\geq Y$，则
\[
\int_\Omega X\,d\mu\geq\int_\Omega Y\,d\mu.
\]
\item 若 $Y$ 是非负简单函数，$\{X_n\}$ 是逐点单调不减的非负简单函数列，且处处有 $\lim_{n\to\infty}X_n\geq Y$，则
\[
\lim_{n\to\infty}\int_\Omega X_n\,d\mu
\geq\int_\Omega Y\,d\mu.
\]
\end{enumerate}
\end{proposition}

\begin{proof}
\begin{enumerate}
\item 以 $A$ 和 $\Omega\setminus A$ 为划分，由定义及零系数约定，
\[
\int_\Omega I_A\,d\mu
=1\cdot\mu(A)+0\cdot\mu(\Omega\setminus A)=\mu(A).
\]
\item 写成 $X=\sum_{i=1}^N a_iI_{A_i}$，其中 $a_i\geq0$，则
\[
\int_\Omega X\,d\mu=\sum_{i=1}^N a_i\mu(A_i)\geq0.
\]
\item 若 $a=0$，两侧均为 $0$。若 $a>0$，沿用 $X$ 的划分，有
\begin{align*}
\int_\Omega aX\,d\mu
&=\sum_{i=1}^N aa_i\mu(A_i)\\
&=a\sum_{i=1}^N a_i\mu(A_i)
=a\int_\Omega X\,d\mu.
\end{align*}
所有项非负，等式也适用于积分为 $+\infty$ 的情形。
\item 设
\[
X=\sum_{i=1}^N a_iI_{A_i},\qquad
Y=\sum_{j=1}^M b_jI_{B_j}.
\]
在共同细化 $\{A_i\cap B_j\}$ 上，
\[
X+Y=\sum_{i=1}^N\sum_{j=1}^M(a_i+b_j)I_{A_i\cap B_j}.
\]
因此
\begin{align*}
\int_\Omega(X+Y)\,d\mu
&=\sum_{i=1}^N\sum_{j=1}^M(a_i+b_j)\mu(A_i\cap B_j)\\
&=\sum_{i=1}^N a_i\sum_{j=1}^M\mu(A_i\cap B_j)
 +\sum_{j=1}^M b_j\sum_{i=1}^N\mu(A_i\cap B_j)\\
&=\sum_{i=1}^N a_i\mu(A_i)+\sum_{j=1}^M b_j\mu(B_j)\\
&=\int_\Omega X\,d\mu+\int_\Omega Y\,d\mu.
\end{align*}
这里对每个 $i$，集合 $\{A_i\cap B_j\}_{j=1}^M$ 划分 $A_i$；对每个 $j$，集合 $\{A_i\cap B_j\}_{i=1}^N$ 划分 $B_j$，所以可分别应用有限可加性。
\item 因为 $X-Y$ 是非负简单函数，由 (2)、(4) 得
\[
\int_\Omega X\,d\mu
=\int_\Omega(X-Y)\,d\mu+\int_\Omega Y\,d\mu
\geq\int_\Omega Y\,d\mu.
\]
此处不对积分作减法，即使积分为无穷，推理仍成立。
\item 由 (5)，$\int_\Omega X_n\,d\mu$ 单调不减，故其极限
\[
L:=\lim_{n\to\infty}\int_\Omega X_n\,d\mu
\]
在 $[0,+\infty]$ 中存在。固定 $0<\alpha<1$，令
\[
E_n=\{X_n\geq\alpha Y\}.
\]
由于 $X_n,Y$ 都是可测函数，$E_n\in\mathcal F$；由 $X_n\leq X_{n+1}$，有 $E_n\subseteq E_{n+1}$。

若 $Y(\omega)=0$，则 $X_n(\omega)\geq0$，所以该点属于每个 $E_n$。若 $Y(\omega)>0$，则
\[
\lim_{n\to\infty}X_n(\omega)\geq Y(\omega)>\alpha Y(\omega),
\]
因而充分大的 $n$ 满足 $X_n(\omega)\geq\alpha Y(\omega)$。所以
\[
E_n\uparrow\Omega.
\]
设 $Y=\sum_{j=1}^M b_jI_{B_j}$，其中 $b_j\geq0$。由测度的下连续性，
\[
\mu(B_j\cap E_n)\uparrow\mu(B_j).
\]
而 $X_nI_{E_n}$ 和 $\alpha YI_{E_n}$ 均为非负简单函数，且
\[
X_n\geq X_nI_{E_n}\geq\alpha YI_{E_n}.
\]
由 (3)、(5) 及积分定义，
\begin{align*}
\int_\Omega X_n\,d\mu
&\geq\int_\Omega X_nI_{E_n}\,d\mu\\
&\geq\alpha\int_\Omega YI_{E_n}\,d\mu\\
&=\alpha\sum_{j=1}^M b_j\mu(B_j\cap E_n).
\end{align*}
有限个非负项可以逐项取单调极限；零系数项恒为 $0$。令 $n\to\infty$，得
\[
L\geq\alpha\sum_{j=1}^M b_j\mu(B_j)
=\alpha\int_\Omega Y\,d\mu.
\]
若 $\int_\Omega Y\,d\mu=+\infty$，任取一个 $\alpha>0$ 即知 $L=+\infty$；若该积分有限，再令 $\alpha\uparrow1$，得 $L\geq\int_\Omega Y\,d\mu$。
\end{enumerate}
\end{proof}

\subsection*{非负可测函数的积分}

\begin{mydefinition}{非负可测函数的积分}{integral-nonnegative}
设 $X:\Omega\to[0,+\infty]$ 可测，定义
\[
\int_\Omega X\,d\mu
:=
\sup\left\{
\int_\Omega Y\,d\mu:
Y\text{ 为非负简单函数，且 }Y\leq X
\right\}.
\]
其中 $Y\leq X$ 指对每个 $\omega\in\Omega$ 都成立。该值称为 $X$ 关于 $\mu$ 的\textbf{积分}，也写作 $\int_\Omega X(\omega)\,\mu(d\omega)$。
\end{mydefinition}

\textbf{注：}\quad 非负简单函数逼近定理给出了满足 $Y\leq X$ 的简单函数；零函数本身也满足此条件。因此，上述取上确界的集合非空，并且包含于 $[0,+\infty]$，积分总能在 $[0,+\infty]$ 中定义。

\begin{proposition}{非负可测函数积分的性质}{integral-nonnegative-properties}
在测度空间 $(\Omega,\mathcal F,\mu)$ 上，有下列结论。
\begin{enumerate}
\item 若 $X$ 为非负简单函数，则以上两种积分定义给出相同的值。
\item 若 $X_n$ 为非负简单函数，且 $X_n\uparrow X$ 逐点成立，则
\[
\lim_{n\to\infty}\int_\Omega X_n\,d\mu=\int_\Omega X\,d\mu.
\]
\item 若 $X$ 为非负可测函数，则 $\int_\Omega X\,d\mu\geq0$。
\item 若 $X$ 为非负可测函数，$a\geq0$ 为实常数，则
\[
\int_\Omega aX\,d\mu=a\int_\Omega X\,d\mu.
\]
\item 若 $X,Y$ 为非负可测函数，则
\[
\int_\Omega(X+Y)\,d\mu
=\int_\Omega X\,d\mu+\int_\Omega Y\,d\mu.
\]
\item 若非负可测函数 $X,Y$ 处处满足 $X\geq Y$，则
\[
\int_\Omega X\,d\mu\geq\int_\Omega Y\,d\mu.
\]
\end{enumerate}
\end{proposition}

\begin{proof}
\begin{enumerate}
\item 将简单函数积分定义给出的值暂记为 $J_s(X)$，由上确界定义给出的值记为 $J(X)$。因为 $X$ 本身是满足 $X\leq X$ 的非负简单函数，所以
\[
J(X)\geq J_s(X).
\]
对任意满足 $0\leq Y\leq X$ 的简单函数 $Y$，简单函数积分的单调性给出
\[
J_s(Y)\leq J_s(X).
\]
对所有这样的 $Y$ 取上确界，得 $J(X)\leq J_s(X)$，故两种积分相同。
\item 由简单函数积分的单调性，$\int_\Omega X_n\,d\mu$ 单调不减。记其极限为 $L$。由于 $X_n\leq X$，由积分的上确界定义及 (1)，
\[
\int_\Omega X_n\,d\mu\leq\int_\Omega X\,d\mu,
\qquad L\leq\int_\Omega X\,d\mu.
\]
另一方面，对任意满足 $Y\leq X$ 的非负简单函数 $Y$，有
\[
Y\leq X=\lim_{n\to\infty}X_n.
\]
由前述非负简单函数列的比较性质，$L\geq\int_\Omega Y\,d\mu$。对所有这样的 $Y$ 取上确界，得到
\[
L\geq\int_\Omega X\,d\mu.
\]
两向不等式合起来即得结论。
\item 上确界定义中的每个简单函数积分都非负，且零函数给出的值为 $0$，故 $\int_\Omega X\,d\mu\geq0$。
\item 当 $a=0$ 时，约定 $0\cdot(+\infty)=0$，所以两侧均为 $0$。当 $a>0$ 时，由简单函数逼近定理，取 $X_n\uparrow X$，则 $aX_n\uparrow aX$。这也说明 $aX$ 是非负可测函数。利用 (2) 及简单函数积分的齐次性，
\begin{align*}
\int_\Omega aX\,d\mu
&=\lim_{n\to\infty}\int_\Omega aX_n\,d\mu\\
&=\lim_{n\to\infty}a\int_\Omega X_n\,d\mu
=a\int_\Omega X\,d\mu.
\end{align*}
\item 取非负简单函数列 $X_n\uparrow X$、$Y_n\uparrow Y$，则 $X_n+Y_n$ 是非负简单函数，且
\[
X_n+Y_n\uparrow X+Y.
\]
由可测函数列极限的可测性，$X+Y$ 可测。再由 (2) 和简单函数积分的可加性，
\begin{align*}
\int_\Omega(X+Y)\,d\mu
&=\lim_{n\to\infty}\int_\Omega(X_n+Y_n)\,d\mu\\
&=\lim_{n\to\infty}
 \left(\int_\Omega X_n\,d\mu+\int_\Omega Y_n\,d\mu\right)\\
&=\int_\Omega X\,d\mu+\int_\Omega Y\,d\mu.
\end{align*}
最后一步只涉及非负数列的单调极限，即使有积分为 $+\infty$，也不会出现无穷减无穷。
\item 任取满足 $0\leq Z\leq Y$ 的简单函数 $Z$，由 $Y\leq X$，也有 $Z\leq X$。因此
\[
\int_\Omega Z\,d\mu\leq\int_\Omega X\,d\mu.
\]
对所有满足 $0\leq Z\leq Y$ 的简单函数取上确界，便得
\[
\int_\Omega Y\,d\mu\leq\int_\Omega X\,d\mu.
\]
\end{enumerate}
\end{proof}

\subsection*{一般可测函数的积分}

对实值可测函数 $X:\Omega\to\mathbb R$，令
\[
X^+=\max\{X,0\},\qquad X^-=-\min\{X,0\}.
\]
分别称 $X^+$ 和 $X^-$ 为 $X$ 的\textbf{正部}和\textbf{负部}。它们都是非负可测函数，并且逐点有 $X=X^+-X^-$。

\begin{mydefinition}{积分存在与可积性}{integrable-function}
设 $X:\Omega\to\mathbb R$ 可测。若
\[
\min\left\{\int_\Omega X^+\,d\mu,\int_\Omega X^-\,d\mu\right\}<+\infty,
\]
即正部和负部的积分至少有一个有限，则称 $X$ 的\textbf{积分存在}，或称其\textbf{积分有意义}，并定义
\[
\int_\Omega X\,d\mu
:=\int_\Omega X^+\,d\mu-\int_\Omega X^-\,d\mu.
\]
若进一步满足
\[
\max\left\{\int_\Omega X^+\,d\mu,\int_\Omega X^-\,d\mu\right\}<+\infty,
\]
即两个积分都有限，则称 $X$ 关于 $\mu$ \textbf{可积}。
\end{mydefinition}

\textbf{注：}\quad 积分存在时，其值可以是实数、$+\infty$ 或 $-\infty$；可积时积分为有限实数。若正部和负部的积分均为 $+\infty$，则不能用 $+\infty-(+\infty)$ 定义积分。

\subsection*{几乎处处关系与积分的性质}

\begin{mydefinition}{几乎处处成立}{almost-everywhere}
设 $(\Omega,\mathcal F,\mu)$ 为测度空间。若 $A\in\mathcal F$ 满足 $\mu(A^c)=0$，则称 $A$ \textbf{几乎处处成立}，简记为 $A,\ \mathrm{a.e.}$。

若一个条件除去某个可测零测集后处处成立，也称该条件几乎处处成立。例如
\[
X=Y,\ \mathrm{a.e.}\quad\Longleftrightarrow\quad
\mu(\{X\ne Y\})=0,
\]
\[
X\geq Y,\ \mathrm{a.e.}\quad\Longleftrightarrow\quad
\mu(\{X<Y\})=0.
\]
\end{mydefinition}

对 $A\in\mathcal F$，定义
\[
\int_A X\,d\mu:=\int_\Omega I_AX\,d\mu,
\]
只要右侧积分有意义。当 $X$ 的积分存在时，由
\[
(I_AX)^+=I_AX^+,\qquad (I_AX)^-=I_AX^-,
\]
以及 $0\leq I_AX^\pm\leq X^\pm$，可知右侧的两个非负部分至少有一个积分有限，所以 $\displaystyle\int_A X\,d\mu$ 也存在。

\begin{proposition}{积分与几乎处处关系}{integral-almost-everywhere}
以下有符号可测函数取实值，非负可测函数也允许取 $+\infty$。
\begin{enumerate}
\item 若 $X$ 的积分存在，则
\[
\left|\int_\Omega X\,d\mu\right|
\leq\int_\Omega|X|\,d\mu.
\]
\item $X$ 可积当且仅当 $|X|$ 可积。
\item 若 $X$ 的积分存在，且 $A\in\mathcal F$ 满足 $\mu(A)=0$，则
\[
\int_A X\,d\mu=0.
\]
\item 若 $X,Y$ 的积分分别存在，且 $X\geq Y,\ \mathrm{a.e.}$，则
\[
\int_\Omega X\,d\mu\geq\int_\Omega Y\,d\mu.
\]
\item 若 $X=Y,\ \mathrm{a.e.}$，且其中一个的积分存在，则另一个的积分也存在，且
\[
\int_\Omega X\,d\mu=\int_\Omega Y\,d\mu.
\]
\item 若 $X=0,\ \mathrm{a.e.}$，则 $X$ 可积，且 $\displaystyle\int_\Omega X\,d\mu=0$。
\item 若 $X\geq0,\ \mathrm{a.e.}$，且 $\displaystyle\int_\Omega X\,d\mu=0$，则 $X=0,\ \mathrm{a.e.}$。
\end{enumerate}
\end{proposition}

\begin{proof}
\begin{enumerate}
\item 由 $|X|=X^++X^-$ 及非负函数积分的可加性，
\[
\int_\Omega|X|\,d\mu
=\int_\Omega X^+\,d\mu+\int_\Omega X^-\,d\mu.
\]
若右侧为 $+\infty$，所需不等式成立。若右侧有限，则两个积分均有限，由非负实数的差的绝对值不超过其和，得
\[
\left|\int_\Omega X\,d\mu\right|
=\left|\int_\Omega X^+\,d\mu-\int_\Omega X^-\,d\mu\right|
\leq\int_\Omega|X|\,d\mu.
\]
\item 若 $|X|$ 可积，由 $0\leq X^+,X^-\leq|X|$ 及非负积分的单调性，
\[
\max\left\{\int_\Omega X^+\,d\mu,\int_\Omega X^-\,d\mu\right\}
\leq\int_\Omega|X|\,d\mu<+\infty,
\]
所以 $X$ 可积。反过来，若 $X$ 可积，则由上一项的恒等式，
\[
\int_\Omega|X|\,d\mu
=\int_\Omega X^+\,d\mu+\int_\Omega X^-\,d\mu<+\infty,
\]
故 $|X|$ 可积。
\item 先设 $X=\sum_{i=1}^N a_iI_{A_i}$ 为非负简单函数。由 $\mu(A)=0$，有 $\mu(A_i\cap A)=0$，所以
\[
\int_A X\,d\mu
=\sum_{i=1}^N a_i\mu(A_i\cap A)=0.
\]
若 $X$ 为非负可测函数，取非负简单函数 $X_n\uparrow X$，则 $I_AX_n\uparrow I_AX$，由已证的简单函数逼近积分公式，
\[
\int_A X\,d\mu
=\lim_{n\to\infty}\int_A X_n\,d\mu=0.
\]
若 $X$ 为一般实值可测函数，分别对 $X^+,X^-$ 应用上述结论，得到
\[
\int_\Omega I_AX^+\,d\mu
=\int_\Omega I_AX^-\,d\mu=0.
\]
因此限制在 $A$ 上的积分存在，并且等于 $0-0=0$。
\item 先设 $X,Y\geq0$，令 $N=\{X<Y\}$，则 $\mu(N)=0$，并且
\[
I_{N^c}X\geq I_{N^c}Y.
\]
由非负积分的单调性、可加性及 (3)，
\begin{align*}
\int_\Omega X\,d\mu
&=\int_\Omega I_NX\,d\mu+\int_\Omega I_{N^c}X\,d\mu\\
&=\int_\Omega I_{N^c}X\,d\mu\\
&\geq\int_\Omega I_{N^c}Y\,d\mu
=\int_\Omega Y\,d\mu.
\end{align*}

对一般实值函数，先比较正部和负部。若 $X^+<Y^+$，则 $Y>0$；在 $X\leq0$ 时有 $X<Y$，在 $X>0$ 时也有 $X=X^+<Y^+=Y$，所以
\[
\{X^+<Y^+\}\subseteq\{X<Y\}.
\]
若 $X^->Y^-$，则 $X<0$；在 $Y\geq0$ 时有 $X<Y$，在 $Y<0$ 时有 $X=-X^-<-Y^-=Y$，所以
\[
\{X^->Y^-\}\subseteq\{X<Y\}.
\]
故 $X^+\geq Y^+,\ \mathrm{a.e.}$，$X^-\leq Y^-,\ \mathrm{a.e.}$。应用已证的非负情形，得
\[
\int_\Omega X^+\,d\mu\geq\int_\Omega Y^+\,d\mu,\qquad
\int_\Omega X^-\,d\mu\leq\int_\Omega Y^-\,d\mu.
\]
若 $\displaystyle\int_\Omega X\,d\mu=+\infty$ 或 $\displaystyle\int_\Omega Y\,d\mu=-\infty$，结论直接成立。其余情形下，由积分存在的定义，$\displaystyle\int_\Omega X^+\,d\mu$ 与 $\displaystyle\int_\Omega Y^-\,d\mu$ 均有限；上面两式又保证其余两个积分有限。因此可以相减，得到
\[
\int_\Omega X\,d\mu
=\int_\Omega X^+\,d\mu-\int_\Omega X^-\,d\mu
\geq\int_\Omega Y^+\,d\mu-\int_\Omega Y^-\,d\mu
=\int_\Omega Y\,d\mu.
\]
\item 先对非负 $X,Y$ 分别应用 (4) 的两个方向，即得积分相等，包括积分为 $+\infty$ 的情形。

对一般实值函数，由 $X=Y,\ \mathrm{a.e.}$，有 $X^+=Y^+,\ \mathrm{a.e.}$ 和 $X^-=Y^-,\ \mathrm{a.e.}$。由非负情形，
\[
\int_\Omega X^+\,d\mu=\int_\Omega Y^+\,d\mu,\qquad
\int_\Omega X^-\,d\mu=\int_\Omega Y^-\,d\mu.
\]
若 $X$ 的积分存在，则左侧至少有一个有限，右侧也如此，所以 $Y$ 的积分存在。将这两式相减，便得所需等式。若先知 $Y$ 的积分存在，交换两者即可。
\item 对 $X$ 与零函数应用 (5)，便有 $\displaystyle\int_\Omega X\,d\mu=0$。同时对非负函数 $|X|$ 与零函数应用 (5)，得 $\displaystyle\int_\Omega|X|\,d\mu=0$，所以由 (2)，$X$ 可积。
\item 对每个正整数 $k$，有
\[
\frac1k I_{\{X>1/k\}}\leq XI_{\{X>1/k\}}\leq X,
\qquad\mathrm{a.e.}
\]
中间函数非负，且积分不超过 $\displaystyle\int_\Omega X^+\,d\mu$。由 $X\geq0,\ \mathrm{a.e.}$，$X^-=0,\ \mathrm{a.e.}$，所以由 (6) 及积分为零的假设，
\[
\int_\Omega X^-\,d\mu=0,\qquad
\int_\Omega X^+\,d\mu=0.
\]
因此中间函数的积分存在。由 (4)，
\[
0\leq\frac1k\mu(\{X>1/k\})
=\int_\Omega\frac1k I_{\{X>1/k\}}\,d\mu
\leq\int_\Omega XI_{\{X>1/k\}}\,d\mu
\leq\int_\Omega X\,d\mu=0.
\]
于是每个 $\{X>1/k\}$ 都为零测集。由
\[
\{X>0\}=\bigcup_{k=1}^{\infty}\{X>1/k\}
\]
及测度的可数次可加性，
\[
\mu(\{X>0\})\leq\sum_{k=1}^{\infty}\mu(\{X>1/k\})=0.
\]
再结合 $\mu(\{X<0\})=0$，得到 $\mu(\{X\ne0\})=0$。
\end{enumerate}
\end{proof}

\subsection*{可积函数积分的线性}

\begin{mytheorem}{积分的线性}{integral-linearity}
设 $X,Y$ 为测度空间 $(\Omega,\mathcal F,\mu)$ 上的可积实值函数。
\begin{enumerate}
\item 对任意实常数 $a$，$aX$ 可积，且
\[
\int_\Omega aX\,d\mu=a\int_\Omega X\,d\mu.
\]
\item $X+Y$ 可积，且
\[
\int_\Omega(X+Y)\,d\mu
=\int_\Omega X\,d\mu+\int_\Omega Y\,d\mu.
\]
\end{enumerate}
\end{mytheorem}

\begin{proof}
\begin{enumerate}
\item 若 $a\geq0$，则
\[
(aX)^+=aX^+,\qquad (aX)^-=aX^-.
\]
由非负积分的齐次性及 $X$ 可积，
\[
\int_\Omega(aX)^+\,d\mu=a\int_\Omega X^+\,d\mu<+\infty,\qquad
\int_\Omega(aX)^-\,d\mu=a\int_\Omega X^-\,d\mu<+\infty.
\]
所以 $aX$ 可积，且
\begin{align*}
\int_\Omega aX\,d\mu
&=a\int_\Omega X^+\,d\mu-a\int_\Omega X^-\,d\mu\\
&=a\int_\Omega X\,d\mu.
\end{align*}
若 $a<0$，正负部交换：
\[
(aX)^+=(-a)X^-,\qquad (aX)^-=(-a)X^+.
\]
由于 $-a>0$，这两个非负函数的积分均有限，故 $aX$ 可积，并且
\begin{align*}
\int_\Omega aX\,d\mu
&=(-a)\int_\Omega X^-\,d\mu-(-a)\int_\Omega X^+\,d\mu\\
&=a\int_\Omega X\,d\mu.
\end{align*}
\item 先验证一种非负分解的积分公式。若可积函数 $Z$ 写成 $Z=U-V$，其中 $U,V$ 均非负可积，则
\[
U+Z^-=V+Z^+.
\]
由非负积分的可加性，
\[
\int_\Omega U\,d\mu+\int_\Omega Z^-\,d\mu
=\int_\Omega V\,d\mu+\int_\Omega Z^+\,d\mu.
\]
所有积分均有限，可以移项，得到
\[
\int_\Omega Z\,d\mu
=\int_\Omega Z^+\,d\mu-\int_\Omega Z^-\,d\mu
=\int_\Omega U\,d\mu-\int_\Omega V\,d\mu.
\]

现令 $U=X^++Y^+$、$V=X^-+Y^-$，则 $X+Y=U-V$。由于 $X,Y$ 可积，由非负积分的可加性，
\[
\int_\Omega U\,d\mu
=\int_\Omega X^+\,d\mu+\int_\Omega Y^+\,d\mu<+\infty,
\]
\[
\int_\Omega V\,d\mu
=\int_\Omega X^-\,d\mu+\int_\Omega Y^-\,d\mu<+\infty.
\]
又由 $-V\leq X+Y\leq U$，有
\[
(X+Y)^+\leq U,\qquad (X+Y)^-\leq V.
\]
利用非负积分的单调性，两个正负部的积分均有限，所以 $X+Y$ 可积。应用刚验证的非负分解公式，
\begin{align*}
\int_\Omega(X+Y)\,d\mu
&=\int_\Omega U\,d\mu-\int_\Omega V\,d\mu\\
&=\left(\int_\Omega X^+\,d\mu+\int_\Omega Y^+\,d\mu\right)
 -\left(\int_\Omega X^-\,d\mu+\int_\Omega Y^-\,d\mu\right)\\
&=\left(\int_\Omega X^+\,d\mu-\int_\Omega X^-\,d\mu\right)
 +\left(\int_\Omega Y^+\,d\mu-\int_\Omega Y^-\,d\mu\right)\\
&=\int_\Omega X\,d\mu+\int_\Omega Y\,d\mu.
\end{align*}
\end{enumerate}
\end{proof}

\begin{mytheorem}{由集合上的积分判别大小与相等}{integral-order-determination}
设 $X,Y$ 为测度空间 $(\Omega,\mathcal F,\mu)$ 上的可积实值函数。
\begin{enumerate}
\item 若对每个 $A\in\mathcal F$ 都有
\[
\int_A X\,d\mu\geq\int_A Y\,d\mu,
\]
则 $X\geq Y,\ \mathrm{a.e.}$。
\item 若对每个 $A\in\mathcal F$ 都有
\[
\int_A X\,d\mu=\int_A Y\,d\mu,
\]
则 $X=Y,\ \mathrm{a.e.}$。
\end{enumerate}
\end{mytheorem}

\begin{proof}
\begin{enumerate}
\item 令 $A=\{X<Y\}\in\mathcal F$。在 $A$ 上 $X<Y$，在 $A^c$ 上两个示性函数乘积均为 $0$，所以 $I_AX\leq I_AY$ 处处成立。

由 $|I_AX|\leq|X|$ 和 $|I_AY|\leq|Y|$，这两个限制函数的绝对值积分均有限，因而都是可积函数。由积分的单调性及题设，
\[
\int_A X\,d\mu\leq\int_A Y\,d\mu
\leq\int_A X\,d\mu,
\]
故两个积分相等。再由积分的线性，
\[
\int_\Omega I_A(Y-X)\,d\mu
=\int_A Y\,d\mu-\int_A X\,d\mu=0.
\]
函数 $I_A(Y-X)$ 非负，由非负函数积分为零的判别，它几乎处处为零。而在 $A$ 上它严格大于 $0$，在 $A^c$ 上等于 $0$，所以
\[
A=\{I_A(Y-X)\ne0\},\qquad\mu(A)=0.
\]
即 $X\geq Y,\ \mathrm{a.e.}$。
\item 等式蕴含两个方向的不等式，分别应用 (1)，得 $X\geq Y,\ \mathrm{a.e.}$ 和 $Y\geq X,\ \mathrm{a.e.}$。因此
\[
\mu(\{X\ne Y\})
\leq\mu(\{X<Y\})+\mu(\{Y<X\})=0,
\]
即 $X=Y,\ \mathrm{a.e.}$。
\end{enumerate}
\end{proof}

\subsection*{几乎处处收敛与单调收敛定理}

\begin{mydefinition}{几乎处处收敛}{convergence-almost-everywhere}
设 $X_n$（$n\geq1$）和 $X$ 为同一测度空间上的可测函数。若存在 $A\in\mathcal F$，使 $\mu(A^c)=0$，且对每个 $\omega\in A$，
\[
\lim_{n\to\infty}X_n(\omega)=X(\omega),
\]
则称 $X_n$ \textbf{几乎处处收敛}于 $X$，记为
\[
X_n\xrightarrow{\mathrm{a.e.}}X.
\]
\end{mydefinition}

\textbf{注：}\quad 几乎处处收敛允许在一个固定的零测集上不收敛；去掉该零测集后必须逐点收敛。逐点收敛也蕴含几乎处处收敛。

\begin{mytheorem}{Levi 单调收敛定理}{levi-monotone-convergence}
设 $X_n$（$n\geq1$）和 $X$ 为测度空间 $(\Omega,\mathcal F,\mu)$ 上的非负可测函数，允许取 $+\infty$。若 $X_n\leq X_{n+1}$ 处处成立，且 $X_n\xrightarrow{\mathrm{a.e.}}X$，则
\[
\lim_{n\to\infty}\int_\Omega X_n\,d\mu
=\int_\Omega\lim_{n\to\infty}X_n\,d\mu
=\int_\Omega X\,d\mu.
\]
上述等式允许各积分为 $+\infty$。
\end{mytheorem}

% 材料未提供 Levi 定理证明；以下用积分的上确界定义与 alpha 截集补证，
% 沿用前述简单函数列比较性质的方法，不引用本定理本身或后续收敛定理。
\begin{proof}
取 $A\in\mathcal F$，使 $\mu(A^c)=0$，且在 $A$ 上 $X_n\to X$。将所有 $X_n$ 和 $X$ 在 $A^c$ 上同时改为 $0$，得到的函数仍非负可测、仍逐点单调，且处处收敛。由几乎处处相等不改变非负积分，所有积分值均不变。因此先证明处处有 $X_n\uparrow X$ 的情形。

由非负积分的单调性，$\displaystyle\int_\Omega X_n\,d\mu$ 单调不减，故其极限 $L\in[0,+\infty]$ 存在。又由 $X_n\leq X$，
\[
L\leq\int_\Omega X\,d\mu.
\]
任取非负简单函数 $Y=\sum_{j=1}^M b_jI_{B_j}\leq X$，固定 $0<\alpha<1$，令
\[
E_n=\{X_n\geq\alpha Y\}.
\]
由于 $X_n$ 单调，$E_n$ 递增。若 $Y(\omega)=0$，该点属于每个 $E_n$；若 $Y(\omega)>0$，由 $X_n(\omega)\uparrow X(\omega)\geq Y(\omega)>\alpha Y(\omega)$，该点最终属于 $E_n$。因此 $E_n\uparrow\Omega$。

非负简单函数 $\alpha YI_{E_n}$ 不超过 $X_n$。由积分的上确界定义及简单函数积分公式，
\[
\int_\Omega X_n\,d\mu
\geq\alpha\sum_{j=1}^M b_j\mu(B_j\cap E_n).
\]
由测度的下连续性，$\mu(B_j\cap E_n)\uparrow\mu(B_j)$。有限个非负项可逐项取极限，零系数项按约定恒为 $0$，于是
\[
L\geq\alpha\sum_{j=1}^M b_j\mu(B_j)
=\alpha\int_\Omega Y\,d\mu.
\]
若 $\displaystyle\int_\Omega Y\,d\mu=+\infty$，此式已给出 $L=+\infty$；若该积分有限，令 $\alpha\uparrow1$，得 $L\geq\displaystyle\int_\Omega Y\,d\mu$。再对所有满足 $Y\leq X$ 的非负简单函数取上确界，得到
\[
L\geq\int_\Omega X\,d\mu.
\]
结合反向不等式，证明了处处收敛时的结论，也就证明了原来的几乎处处收敛情形。

原函数列因非负且单调，在每点都有扩展实值极限 $Z=\lim_{n\to\infty}X_n$，且 $Z$ 可测。由假设，$Z=X,\ \mathrm{a.e.}$，所以 $\displaystyle\int_\Omega Z\,d\mu=\int_\Omega X\,d\mu$。这给出定理中间的等号。
\end{proof}
